% Implementation Results -- handbook section 3

\subsection{Training behaviour}
\label{subsec:results-1}
% Required: Loss and metric curves; selected checkpoint; overfitting or underfitting; training time; implementation issues and fixes.


\subsection{Quantitative results}
\label{subsec:results-2}
% Required: Primary and secondary metrics; parameter counts; training and inference time; configurations; mean and standard deviation over multiple runs where feasible.
% Also required: Section 20.1, segmentation: mIoU and Dice; instance segmentation also reports mask AP. Section 21: report hardware, software versions, seeds, checkpoint rules, training time, and inference cost.


\subsection{Qualitative results}
\label{subsec:results-3}
% Required: Correct predictions; difficult but correct cases; failures; predictions against ground truth; task-appropriate visualizations.


\subsection{Comparison and discussion}
\label{subsec:results-4}
% Required: Which method is better and by how much; where differences appear; causes in data, architecture, loss or training; accuracy, speed and parameter trade-offs; whether hypotheses are supported; conclusions that cannot yet be drawn. Listing metrics and declaring the highest score best is not sufficient.
% Also required: Section 20: compute-cost analysis is mandatory alongside the quantitative comparison.


\subsection{Error analysis}
\label{subsec:results-5}
% Required: Per error group: description; frequency or rate if possible; examples; hypothesized cause; improvement ideas.


\subsection{Ablation study}
\label{subsec:results-6}
% Required: Mandatory for Assignments 2 and 3; recommended for Assignment 1 when comparing representations or regularizers.
% Also required: Section 21: at least one ablation is mandatory for Assignment 2.


\subsection{Limitations and conclusion}
\label{subsec:results-7}
% Required: Data, method and experimental limits; generalization; future work; short answer to the research question.
